% GENERATED FILE. Edit canonical lessons, then run npm run generate:books.
% canonical-source-hash: fnv1a64:9afbb9d26a326adc
% canonical-lessons: TE-C141-tiyani, TE-C141-cedaina, TE-C141-karamaina, TE-C141-dhairyam, TE-C141-somari

\chapter{Sweet, Bitter-tasting, Spicy, Courage, Lazy}
\label{ch:te-sweet-bitter-tasting-spicy-courage-lazy}
\hlchaptermodality{141}

\begin{chapteropening}
\textbf{By the end of this chapter:} \emph{I can say sweet, bitter-tasting, spicy, courage and lazy.}

Complete the last lesson of chapter 141: i can say sweet, bitter-tasting, spicy, courage and lazy.
\end{chapteropening}

\section[tīyani]{\te{తీయని} (tīyani) — sweet (of taste)}

\label{lesson:TE-C141-tiyani}



\begin{quote}
Before the new one: say the Telugu for golden, then the Telugu for spoiled.
\end{quote}

\subsection*{You'll want to know: \te{తీయని}}
\textbf{\te{తీయని}} — \emph{tīyani} — \textquotedblleft{}sweet (of taste)\textquotedblright{}.

\begin{grammarlens}[title={What you've built}]
One more word for this chapter.
\end{grammarlens}

\subsection*{Your turn}
\begin{itemize}
\raggedright
  \item \emph{Say it:} \emph{tīyani}
  \item \emph{Say it:} \emph{tīyani}, once more
  \item \emph{Say it:} say \emph{tīyani}
  \item \emph{Recall:} say the Telugu for golden, then the Telugu for spoiled, then say \emph{tīyani} again
\end{itemize}

\begin{tcolorbox}[breakable,colback=teal!4,colframe=teal!35!black,title={Before you move on}]
What does \te{తీయని} mean? (\textquotedblleft{}Sweet (of taste)\textquotedblright{}.) Say it once more.
\end{tcolorbox}
\section[cēdaina]{\te{చేదైన} (cēdaina) — bitter-tasting}

\label{lesson:TE-C141-cedaina}



\begin{quote}
Before the new one: say the Telugu for spoiled, then the Telugu for sweet.
\end{quote}

\subsection*{You'll want to know: \te{చేదైన}}
\textbf{\te{చేదైన}} — \emph{cēdaina} — \textquotedblleft{}bitter-tasting\textquotedblright{}.

\begin{grammarlens}[title={What you've built}]
One more word for this chapter.
\end{grammarlens}

\subsection*{Your turn}
\begin{itemize}
\raggedright
  \item \emph{Say it:} \emph{cēdaina}
  \item \emph{Say it:} \emph{cēdaina}, once more
  \item \emph{Say it:} say \emph{cēdaina}
  \item \emph{Recall:} say the Telugu for spoiled, then the Telugu for sweet, then say \emph{cēdaina} again
\end{itemize}

\begin{tcolorbox}[breakable,colback=teal!4,colframe=teal!35!black,title={Before you move on}]
What does \te{చేదైన} mean? (\textquotedblleft{}Bitter-tasting\textquotedblright{}.) Say it once more.
\end{tcolorbox}
\section[kāramaina]{\te{కారమైన} (kāramaina) — spicy}

\label{lesson:TE-C141-karamaina}



\begin{quote}
Before the new one: say the Telugu for sweet, then the Telugu for bitter-tasting.
\end{quote}

\subsection*{You'll want to know: \te{కారమైన}}
\textbf{\te{కారమైన}} — \emph{kāramaina} — \textquotedblleft{}spicy\textquotedblright{}.

\begin{grammarlens}[title={What you've built}]
One more word for this chapter.
\end{grammarlens}

\subsection*{Your turn}
\begin{itemize}
\raggedright
  \item \emph{Say it:} \emph{kāramaina}
  \item \emph{Say it:} \emph{kāramaina}, once more
  \item \emph{Say it:} say \emph{kāramaina}
  \item \emph{Recall:} say the Telugu for sweet, then the Telugu for bitter-tasting, then say \emph{kāramaina} again
\end{itemize}

\begin{tcolorbox}[breakable,colback=teal!4,colframe=teal!35!black,title={Before you move on}]
What does \te{కారమైన} mean? (\textquotedblleft{}Spicy\textquotedblright{}.) Say it once more.
\end{tcolorbox}
\section[dhairyaṁ]{\te{ధైర్యం} (dhairyaṁ) — courage; brave}

\label{lesson:TE-C141-dhairyam}



\begin{quote}
Before the new one: say the Telugu for bitter-tasting, then the Telugu for spicy.
\end{quote}

\subsection*{You'll want to know: \te{ధైర్యం}}
\textbf{\te{ధైర్యం}} — \emph{dhairyaṁ} — \textquotedblleft{}courage; brave\textquotedblright{}.

\begin{grammarlens}[title={What you've built}]
One more word for this chapter.
\end{grammarlens}

\subsection*{Your turn}
\begin{itemize}
\raggedright
  \item \emph{Say it:} \emph{dhairyaṁ}
  \item \emph{Say it:} \emph{dhairyaṁ}, once more
  \item \emph{Say it:} say \emph{dhairyaṁ}
  \item \emph{Recall:} say the Telugu for bitter-tasting, then the Telugu for spicy, then say \emph{dhairyaṁ} again
\end{itemize}

\begin{tcolorbox}[breakable,colback=teal!4,colframe=teal!35!black,title={Before you move on}]
What does \te{ధైర్యం} mean? (\textquotedblleft{}Courage; brave\textquotedblright{}.) Say it once more.
\end{tcolorbox}
\section[sōmari]{\te{సోమరి} (sōmari) — lazy}

\label{lesson:TE-C141-somari}



\begin{quote}
Before the new one: say the Telugu for spicy, then the Telugu for courage; brave.
\end{quote}

\subsection*{You'll want to know: \te{సోమరి}}
\textbf{\te{సోమరి}} — \emph{sōmari} — \textquotedblleft{}lazy\textquotedblright{}.

\begin{grammarlens}[title={What you've built}]
One more word for this chapter.
\end{grammarlens}

\subsection*{Your turn}
\begin{itemize}
\raggedright
  \item \emph{Say it:} \emph{sōmari}
  \item \emph{Say it:} \emph{sōmari}, once more
  \item \emph{Say it:} say \emph{sōmari}
  \item \emph{Recall:} say the Telugu for spicy, then the Telugu for courage; brave, then say \emph{sōmari} again
\end{itemize}

\begin{tcolorbox}[breakable,colback=teal!4,colframe=teal!35!black,title={Before you move on}]
What does \te{సోమరి} mean? (\textquotedblleft{}Lazy\textquotedblright{}.) Say it once more.
\end{tcolorbox}
